% !TeX root = ../thesis.tex
\section{Questions with more than two SQL attempts}
\label{app:extended-sql-trajectories}

\hyperref[tab:extended-sql-trajectories]{Table~\ref*{tab:extended-sql-trajectories}}
lists the standalone questions with more than two recorded SQL attempts
in either Nova run. Counts were taken from the saved case-level evaluations.
Each count includes the first attempt and all subsequent attempts.
A dash means that the question did not exceed two attempts in that run;
it does not indicate missing data. Neither Haiku run contained a question
with more than two attempts.

\begin{table}[H]
\centering
{\small
\setlength{\tabcolsep}{8pt}
\renewcommand{\arraystretch}{1.15}
\begin{tabularx}{\textwidth}{@{}>{\raggedright\arraybackslash}X
  *{2}{>{\centering\arraybackslash}p{0.22\textwidth}}@{}}
\toprule
\textbf{Question} & \textbf{Nova 1 attempts} & \textbf{Nova 2 attempts} \\
\midrule
CB200-M039 & 5 & 28 \\
CB200-M051 & -- & 24 \\
CB200-H003 & -- & 29 \\
CB200-H004 & 23 & -- \\
CB200-H007 & -- & 3 \\
CB200-H010 & 4 & -- \\
CB200-H018 & 5 & 10 \\
CB200-H022 & -- & 22 \\
CB200-H024 & 12 & 23 \\
CB200-H025 & -- & 3 \\
CB200-H026 & -- & 3 \\
CB200-H037 & -- & 3 \\
CB200-H038 & 32 & 25 \\
CB200-H049 & -- & 4 \\
\midrule
\textbf{Questions exceeding two attempts} & \textbf{6} & \textbf{12} \\
\textbf{Attempts on those questions} & \textbf{81} & \textbf{177} \\
\bottomrule
\end{tabularx}
}
\caption{Standalone questions exceeding two SQL attempts (Source: Author, 2026).}
\label{tab:extended-sql-trajectories}
\end{table}

These are the cases behind the longer sequences discussed in
\hyperref[sec:results-error-responses]{Section~\ref*{sec:results-error-responses}}.
They include both recovered and unsuccessful outcomes; the attempt count
alone does not describe whether the final result was correct.
